\section{AION Flow Threat Model, Adversarial Security and Dependency Review}
\label{sec:aion-flow-threat-model}

This section records completion of the formal review required by \texttt{AFLOW-1107}. It is
intended to be appended to the existing AION and Pilot technical document. Completion of the
review does not mean that AION Flow is penetration-tested or ready for public enterprise
exposure. The review identifies the remaining release gates explicitly.

\subsection{Security objective and protected assets}

AION Flow allows a customer to assemble replaceable models, evidence sources, harnesses, compute
targets and business capabilities. None of those components may acquire authority merely by being
placed on the canvas. The customer-controlled AION brain remains the authority, memory, policy,
verification and receipt boundary.

Protected assets include the Business Map, private memory, files, Visual Vault credentials,
provider tokens, endpoint details, canonical membership, workspace boundaries, exact approvals,
workflow revisions, execution receipts, learning evidence, residency rules, compute budgets,
publisher identities and update packages.

The desktop canvas, phone interface, browser request, workflow JSON, encrypted import, template,
model output, retrieved evidence, provider response and connector payload are all treated as
untrusted. A model response is never accepted as proof of identity, authority, execution or
success.

\subsection{Principal threats and controls}

The analysis covers tenant and role forgery, oversized or malformed graphs, revision races,
encrypted import abuse, embedded credentials, template substitution, approval replay, prompt and
model injection, data exfiltration, false execution claims and compromised dependencies.

Production graph qualification now applies hard limits to nesting depth, item count, scalar size,
encoded size, node count and edge count before canonical hashing or persistence. Node and edge
records must have valid types and safe scoped identifiers. Revision commits use both optimistic
base-revision checks and an inter-process exclusive lock before atomic replacement, preventing two
local writers from both accepting the same revision.

Encrypted transfer accepts only the declared AION Flow schema, scrypt profile and AES-256-GCM
cipher. It validates Base64 strictly, bounds salt, nonce, ciphertext and plaintext sizes, requires a
minimum password length, authenticates the header and compares the recovered graph hash without a
content-dependent early exit. Import restores neither permissions nor execution authority.

Signed templates now reject recursively nested password, token, authorization, private-key and
API-key fields in addition to literal vault references and redaction markers. Customer bindings
remain outside the signed package as opaque Visual Vault references. Immutable versions,
publisher signatures, pinning, rollback, quarantine and revocation remain enforced.

\subsection{Adversarial verification}

Automated tests prove fail-closed handling of malformed node and edge records, invalid graph
collections, unsafe identifiers, excessive nesting, cyclic structures, oversized scalar content,
dangling edges, duplicate identities, stale revisions, cross-workspace access and forged role
claims. Transfer tests reject unsupported encryption headers, invalid Base64, oversized ciphertext
and weak import passwords. Template tests reject secret-shaped fields at arbitrary graph depth.

These tests supplement rather than replace independent penetration testing. Provider-specific
prompt-injection corpora, distributed concurrency, connector receipt falsification and sustained
resource-exhaustion testing remain part of production qualification.

\subsection{Dependency review and remediation}

The desktop production and build dependency audit reports no known advisories after upgrading
\texttt{electron-builder} from version 24.13.3 to 26.15.3 and regenerating its lock. The minimal
Python product lock also reports no known advisories after upgrading FastAPI, Starlette, Click,
IDNA, \texttt{python-dotenv} and \texttt{python-multipart}.

The wider backend manifest was upgraded across the principal network, cryptography, HTTP, image,
PDF and rendering packages. Two upstream areas do not yet have a complete patched route within the
current dependency family: the ECDSA implementation used transitively by legacy JWT support and a
WeasyPrint advisory. They must be isolated from untrusted verification or rendering inputs, or
replaced, before an exposed enterprise release.

The all-purpose research environment is expressly not the production bill of materials. It also
contains large machine-learning and document toolchains with independent advisories. Customer
packages must install a minimal hashed product lock plus selected feature extras; they must never
clone the developer environment.

\subsection{Release gates}

The review establishes five mandatory gates. First, the current API remains loopback-only and must
not be exposed directly to a VPN, public listener or customer network. Second, remote or multi-user
consequential execution remains disabled until preparation, approval, execution, run control,
evaluation and template lifecycle endpoints all bind a signed device session to canonical
organisation authority and persisted customer policy. Third, the development fallback encryption
key must be replaced during installation and production startup must fail when a secure key is
absent. Fourth, release engineering must generate a signed software bill of materials and minimal
hashed lock while isolating or replacing the remaining ECDSA and rendering paths. Fifth, an
independent penetration, secrets, licence and update-signature review is required before enterprise
field deployment. Sixth, the macOS package must be signed and notarized with a suitable Developer
ID Application identity; the local qualification package built successfully but remained unsigned.

Accordingly, \texttt{AFLOW-1107} is complete as a formal threat and dependency review with
implemented input hardening. It is not a security certification. The later security-remediation and
load-qualification section records closure of the signed-session, persisted-policy, customer-key,
minimal-lock, signed-SBOM and sustained-load engineering gates. Independent accessibility,
regulatory, penetration and representative field qualification under \texttt{AFLOW-1109} remains
open and cannot be waived by a model, template or canvas policy.
